\chapter{推理优化}
\label{lab:A50_inference_optimization}

\noindent\textbf{配套文件：}\path{notebook/A50_inference_optimization.ipynb}\par
\noindent\textbf{教材关联：}\bookxrefshort{ch:long-context}、\bookxrefshort{ch:incremental-decoding-cache}、\bookxrefshort{ch:inference-algorithms}、\bookxrefshort{ch:quantization}、\bookxrefshort{ch:compute-infrastructure}、\bookxrefshort{ch:batching-scheduling}、\bookxrefshort{ch:serving-architecture}。

\section*{实验目标}
把数值等价性、内存节省和时延收益分别验证。

\section*{准备及定位}
先阅读本册实验总则，确认Notebook中的依赖与资产单元。原理计算、库接口迁移和真实模型训练可能具有不同资源要求，应分别记录设备、版本及运行范围。

源码定位可从下列函数开始；应结合前置数据与配置单元阅读。
\begin{itemize}
\item \texttt{my\_format\_prompt}
\item \texttt{my\_encode\_sft\_record}
\item \texttt{my\_move\_batch}
\end{itemize}

\section*{操作及观察}
\begin{enumerate}
\item 固定输入与模型，比较朴素生成和KV缓存生成的词元及输出差异。
\item 跟踪缓存长度、头数和内存估算，检查左padding、分桶及批处理状态。
\item 比较SDPA、量化和编译路径时记录设备、预热、同步、输入形状及版本。
\item 分别报告TTFT、TPOT和吞吐，审查连续批处理、推测解码及并行部署的适用前提。
\end{enumerate}

\section*{结果判读及提交}
提交缓存轨迹、等价误差与分阶段时间。状态转换演算不代表真实服务压测；前缀缓存还需考虑会话身份和版本隔离。

提交时附上所执行单元范围、环境与制品身份、原始结果及失败说明。将未运行项目明确列出，不能用Notebook中已有输出替代本次执行证据。

相关方法的原始定义与适用前提见\cite{dao2022flashattention,leviathan2023speculative}；本实验的运行结果仍须按实际配置单独验证。
